\section{Lecture November 16}

\begin{defbox}
    Let $R$ a commutative ring, $a \in R$. The \bold{evaluation map at $a$} is the function $\epsilon_a \colon R(x) \to R$ defined by $\sum_{i = 0}^n c_ix^i \mapsto \sum_{i = 0}^n c_ia^i$.
\end{defbox}

\begin{propbox}
    The evaluation maps are a surjective homeomorphism.
\end{propbox}

\begin{proof}
    For each $c \in R$, we have $\epsilon_a(c) = c$, so $\epsilon_a$ is surjective. Furthermore, $\epsilon_a(1) = 1$ and it is straightforward to check that $\epsilon_a$ preserves $+$ and $\cdot$. Indeed, for $f := \sum_{i=0}^nc_ix^i$ and $g := \sum_{j = 0}^m d_jx^j$ we have assume $n = m$ by adding zero coefficients, so
    \[ \epsilon_a(f+g) = \sum_{i = 0}^n(c_i + d_i)a^i = \sum_{i = 0}^{n}c_ia^i + \sum_{i = 0}^{n}d_ia^i = \epsilon_a(f)+\epsilon_a(g). \]
    Moreover, we have that
    \[ fg = \sum_{i,j}c_id_jx^{i+j} = \sum_{k = 0}^{2n}\left(\sum_{i + j = k} c_id_j\right)x^k\]
    thus 
    \[ \epsilon_a(fg) = \sum_{k = 0}^{2n}\left(\sum_{i + j = k} c_id_j\right)a^k = \epsilon_a(f)\epsilon_a(g) \qedhere\]
\end{proof}

\subsection{Remainder and Roots}

\begin{defbox}
    A \bold{root} of a polynomial $f \in R[x]$ where $R$ is some commutative ring, is any element $r \in R$ such that $f(r) = 0$. 
\end{defbox}

\begin{propbox}
    Let $R$ be a commutative ring and $a \in R$. Then for each $f \in R[x]$, $a$ is a root of $f$ if and only if $x - a$ is a factor of $f$, i.e $(x - a) \mid f$.
\end{propbox}

\begin{proof}
    The reverse direction is trivial; indeed, $x = a$ automatically gives zero. For the converse direction, we divide $f$ by $(x-a)$ with remainder to obtain $f = g(x-a) + r$ for somem $g \in R[x]$, where $\deg(r) < \deg(x-a) = 1$, so $r \in R$. Now we have that $0 = f(a) = g(a)(a-a) + r = r$, hence $r = 0$ and $f = g(x-a)$, i.e. $(x-a)\mid f$. 
\end{proof}

\begin{defbox}
    Let $R$ be a commutative ring. An element $a \in R$ is called \bold{irreducible} if $a$ is not a unit and whenever $a = uv$, then either $u$ or $v$ is a unit.
\end{defbox}

Note that for a field $F$, the units $F[x]^\times = F^\times$, hence a polynomial in a field is irreducible if it has positive degree and it is not equal to $fg$ where both $f$ and $g$ have positive degree.

Thus, irreducible elements are a generalization of primes in $\ZZ$.

\begin{propbox}
    Let $F$ be a field and $f \in F[x]$ be of degree $n \geq $ and $c$ is its leading root.
    \begin{enumerate}
        \item $f$ factors into a product of $k \leq n$ irreducibles times $c$, i.e. $f = cp_1\dots p_k$ where $k \leq n$ and the $p_i \in F[x]$ are monic irreducible polynomials.
        \item $f$ has at most $n$ roots.
        \item \textbf{Fundamental Theorem of Algebra}. For each $f \in \CC[x]$, there exist not necessarily distinct $r_1, \dots, r_n$, where $n := \deg(f)$ such that $f = c(x-r_i)\dots(x-r_n)$ for some $c \in C$. 
    \end{enumerate}
\end{propbox}

\begin{proof}
    In HW, by induction.
\end{proof}

\subsection{Ideals}

For a ring homeomorphism $\phi \colon R \to S$, its kernel $\ker(\phi)$ is not a subring unless $S$ is trivial, because it contains $1$. However, $\phi$ is a group homeomorphism $(R, +) \to (S, +)$ so $\ker(\phi)$ is an additive subgroup of $R$. Moreover, if $k \in \ker(\phi)$ and $r \in R$< then 
    \[ \phi(rk) = \phi(r)\phi(k) = \phi(r) \cdot 0 = 0 \]
and similarily $\phi(kr) = 0$, so $R \cdot \ker(\phi) = \ker(\phi) = \ker(\phi) \cdot R$. 

\begin{defbox}
    A subset $I$ of a ring $R$ is caleld a \bold{two-sided ideal} if $I$ is a subgroup under $+$ and $R \cdot I = I = I \cdot R$. Note that $R \cdot I \subseteq I$ already implies $R \cdot I = I$ because $1 \in R$ and $1 \cdot I = I$, so $R \cdot I \supseteq I$ is automatic.
\end{defbox}

\begin{defbox}
    Let $R$ be a commutative ring. The \bold{principal ideal generated by} $k \in R$ is 
    \[ (k) := kR = Rk := \{kr : r \in R\} \]
\end{defbox}

This is indeed an ideal, since if $kr_1$, $kr_2 \in kR$, then $kr_1 + kr_2 = k(r_1 + r_2) \in kR$ and $-kr_1 = k(-r_1) \in kR$. 